% Sections 4.1 and 4.2, from lectures/L04/appendix/a_theory.md: one section per numbered section
% of the appendix.

\section{Neural Network Architecture}
\label{c4:sec:architecture}\label{c4:app:a}

\subsection{Overview}
Chapter~\ref{c3:ch:nn} built a \code{dense_layer::Interface} and a placeholder implementation of
it, \code{dense_layer::Stub}. This chapter adds the piece that actually uses them: a small neural
network class, \code{neural_network::Shallow}, completing the structure:

\begin{console}
ml::dense_layer::Interface   (built in Chapter 3)
└── ml::dense_layer::Stub    (built in Chapter 3)

ml::neural_network::Interface   (Interface for neural networks)
└── ml::neural_network::Shallow (Neural network with a single hidden layer)
        ├── ml::dense_layer::Interface& myHiddenLayer
        └── ml::dense_layer::Interface& myOutputLayer
\end{console}

\noindent \code{ml::dense_layer::Stub} continues to serve as a placeholder throughout this chapter.
Once a concrete implementation is created in Chapter~\ref{c5:ch:dense}, the stub is replaced
without needing to change any of the code you write in this one.

\subsection{The network's structure}
\code{Shallow} holds references to two dense layers and connects them during prediction:

\begin{console}
input → myHiddenLayer.feedforward(input)
               ↓
       myHiddenLayer.output()
               ↓
       myOutputLayer.feedforward(...)
               ↓
       myOutputLayer.output()  →  prediction
\end{console}

\noindent \code{Shallow} doesn't own the layers; it receives them as references via the
constructor. This makes it easy to swap out layers without changing the network class, which is
exactly what happens in Chapter~\ref{c5:ch:dense} when the stub is replaced with a real
implementation.

% ------------------------------------------------------------------------------------------
\section{Training Loop}
\label{c4:sec:training}

\subsection{Overview}
You'll implement the method \code{train()} in \code{ml::neural_network::Shallow}. The method
carries out the complete training process, repeated for every training set and epoch:
\begin{itemize}
  \item Feedforward:
  \begin{itemize}
    \item Computes outputs for every node in the network.
    \item The output from the last layer forms the network's prediction.
    \item Must be computed forward (from input to output).
  \end{itemize}
  \item Backpropagate:
  \begin{itemize}
    \item Computes gradients (simplified error) for every node in the network.
    \item Must be computed backward (from output to input).
  \end{itemize}
  \item Optimization:
  \begin{itemize}
    \item Adjusts the trainable parameters (weights and biases) based on the computed error,
      together with the learning rate the network currently holds.
    \item Computation can happen in either direction.
  \end{itemize}
\end{itemize}

\subsection{Structure of the training loop}
The training loop can be summarized as follows:

\begin{console}
Reset both layers' trainable parameters

For each epoch:
    For each training set x:
        1. Feedforward: compute predictions
        2. Backpropagation: compute error
        3. Optimization: update weights and biases

    Every hundredth epoch:
        4. Measure the precision, and stop early or set a new learning
           rate from it
\end{console}

\noindent The order matters:
\begin{itemize}
  \item Backpropagation depends on the result of feedforward.
  \item Optimization depends on the result of backpropagation.
\end{itemize}
The reset at the top is what makes a second call to \code{train()} train a new network rather than
continue the old one, and step~4 is what frees the caller from having to pick a learning rate. Both
are covered in the exercises, in \secref{c4:sec:exercises}.

\subsection{Step 1: Feedforward}
Call \code{predict(myTrainInput[x])}. This performs feedforward through both layers and returns
the output layer's output.

\begin{console}
predict(input)
  → myHiddenLayer.feedforward(input)
  → myOutputLayer.feedforward(myHiddenLayer.output())
\end{console}

\subsection{Step 2: Backpropagation}
Compute the error backward: always start with the output layer.

\textbf{The output layer} compares its output against the reference value:

\begin{console}
myOutputLayer.backpropagate(myTrainOutput[x])
\end{console}

\noindent\textbf{The hidden layer} computes its error from the output layer's error and weights:

\begin{console}
myHiddenLayer.backpropagate(myOutputLayer)
\end{console}

\noindent The order is crucial: the hidden layer's calculation depends on the output layer having
already computed its error.

\subsection{Step 3: Optimization}
Update the weights and biases in both layers using the errors computed in step~2.

\begin{console}
myHiddenLayer.optimize(myTrainInput[x], learningRate)
myOutputLayer.optimize(myHiddenLayer.output(), learningRate)
\end{console}

\noindent The hidden layer's input is the original input. The output layer's input is the hidden
layer's output (the one computed during feedforward). \code{learningRate} is the network's own, a
local variable in \code{train()} rather than an argument, revised in step~4 below.

\subsection{Step 4: Evaluation}
Once every hundredth epoch, measure how close the network is to its training data and act on the
number:

\begin{console}
precision = 1.0 - MAE (over every value of every training set)

precision >= precisionThreshold  ->  stop training, the network is done
otherwise                        ->  updateLearningRate(): raise the rate when
                                     progress has stalled, lower it when the
                                     network is overshooting
\end{console}

\noindent The learning rate is therefore never supplied by the caller. It starts at \code{0.1} and
is revised from one evaluation to the next, exactly as in \code{ml::lin_reg::Adaptive} from
Chapter~\ref{c2:ch:training}, which evaluates every tenth epoch instead: a linear regression model
converges in tens of epochs, a network in thousands.

\subsection{Connection to the theory}
These three steps map directly onto the theory from Chapter~\ref{c3:ch:nn} (see
\secref{c3:sec:nn}):

\begin{booktable}{@{}p{7.6em}p{11em}L@{}}
  \thead{Step} & \thead{Theory (\secref{c3:sec:nn})} & \thead{Code} \\
  \midrule
  Feedforward & $y = \sigma(b + \sum w_i x_i)$ & \code{feedforward(input)} \\
  Backpropagation & $\Delta e = \delta \cdot y'$ & \code{backpropagate(...)} \\
  Optimization & $w = w + \Delta e \cdot L \cdot x$ & \code{optimize(input, learningRate)} \\
\end{booktable}

\subsection{Input validation}
Check that training is possible before the training loop starts. Print an error message and return
\code{false} immediately if any of the following conditions hold:

\begin{booktable}{@{}p{13.5em}L@{}}
  \thead{Condition} & \thead{Explanation} \\
  \midrule
  \code{epochCount == 0} & Training must run for at least one epoch. \\
  \code{precisionThreshold <= 0.0 || precisionThreshold >= 1.0} & Invalid threshold: must be in
    the range \code{(0.0, 1.0)}. A threshold of \code{1.0} or more can never be reached;
    \code{0.0} or less is reached by any network at all. \\
\end{booktable}

There's no learning rate among them: the network sets its own, and the layers check the rate
they're handed on every call to \code{optimize()}. Missing training data isn't among them either:
it never reaches \code{train()}, because the constructor already terminates when there's no
complete training set.
